\subsection{Big Bang For Your Helium Buck --- arXiv:2609.13140}

\textbf{Authors:} Marilena Loverde, Murali M. Saravanan, Zachary J. Weiner

\paragraph{主な主張}
Large Binocular Telescope による原始ヘリウム存在量の直接測定を prior として用いると、standard BBN の理論関係を仮定しなくても、CMB からの $N_{\rm eff}$、バリオン密度、インフレーション・パラメータの制約を BBN-consistent 解析とほぼ同じ精度まで回復できることを示した \cite{Loverde:2026helium}。

\paragraph{新規性}
CMB 解析で通常用いられる BBN による $Y_{\rm p}$ の理論 calibration を高精度の天文学的ヘリウム測定で置き換え、CMB cosmology の頑健性を BBN 仮定から切り離すと同時に、BBN 自体を new physics の独立 probe として利用できる構成を示した点にある。

\paragraph{位置付け}
原始元素存在量と CMB を結ぶ early-universe precision cosmology の研究であり、dark radiation や BBN--recombination 間の放射成分の進化、さらに基本定数の時間変化を伴う ultralight scalar / modulus 模型を検証する際の観測的基準になる。
